\section{Exact edge separation for double ladders}
\label{sec:exact-hsep}

Throughout this section, let \(a,b\ge3\) be odd.  In ladder \(A\), write
\(r_A(p)=A_p^0A_p^1\) for rung \(p\), \(0\le p\le a\), and
\(\ell_A(i,s)=A_i^sA_{i+1}^s\) for rail \(s\) of cell \(i\),
\(0\le i<a\).  Use \(r_B(q)\) and \(\ell_B(j,s)\) analogously.  Denote
the all-rails cycle by \(R\), the cycle turning in cell \(i\) of \(A\) and
cell \(j\) of \(B\) by \(T(i,j)\), and the four exceptional cycles from
Theorem~\ref{thm:cycle-formula} by
\[
 X_{0011},\quad X_{1100},\quad X_{0101},\quad X_{1010}.
\]
The subscripts record connector words in the order of \cref{eq:connectors}.

\subsection{A symbolic packing lower bound}

We first construct ordered edge requirements whose cycle supports are
pairwise disjoint.  Four requirements isolate the four exceptional cycles.
One concrete choice, anchored on ladder \(A\), is
\begin{equation}
\label{eq:exception-requirements}
 (r_A(0),r_A(1)),\quad (r_A(a),r_A(a-1)),\quad
 (\ell_A(0,0),\ell_A(0,1)),\quad
 (\ell_A(0,1),\ell_A(0,0)).
\end{equation}
The incidence recurrence in Theorem~\ref{thm:cycle-formula} shows directly that the
four supports are distinct singleton sets, one for each \(X\)-cycle.

Next isolate every boundary turn.  For \(0\le j<b\), use
\begin{equation}
\label{eq:horizontal-boundary-requirements}
 (r_A(0),\ell_B(j,j\bmod2)),\qquad
 (r_A(a),\ell_B(j,1-j\bmod2)).
\end{equation}
They have singleton supports \(\{T(0,j)\}\) and
\(\{T(a-1,j)\}\), respectively.  For \(1\le i\le a-2\), the symmetric
requirements
\begin{equation}
\label{eq:vertical-boundary-requirements}
 (r_B(0),\ell_A(i,i\bmod2)),\qquad
 (r_B(b),\ell_A(i,1-i\bmod2))
\end{equation}
have supports \(\{T(i,0)\}\) and \(\{T(i,b-1)\}\).  The parity choice in
\cref{eq:horizontal-boundary-requirements,eq:vertical-boundary-requirements}
ensures that an exceptional cycle containing the endpoint rung also contains
the excluded rail, so no exceptional cycle enters a boundary support.

It remains to handle the interior turn grid.  For \(0<p<a\),
Theorem~\ref{thm:cycle-formula} gives
\begin{equation}
\label{eq:domino-support-a}
 \operatorname{supp}(r_A(p),\ell_B(j,s))
   =\{T(p-1,j),T(p,j)\},
\end{equation}
and, symmetrically, for \(0<q<b\),
\begin{equation}
\label{eq:domino-support-b}
 \operatorname{supp}(r_B(q),\ell_A(i,s))
   =\{T(i,q-1),T(i,q)\}.
\end{equation}
Here either rail side \(s\in\{0,1\}\) may be used.  These requirements are
dominoes in the \((a-2)\times(b-2)\) grid of interior turns.

Delete its far corner \(T(a-2,b-2)\), which belongs to the majority parity
class.  In each interior \(B\)-column, pair the \(A\)-rows
\(1\)--\(2\), \(3\)--\(4\), \(\ldots\), \((a-4)\)--\((a-3)\).
In the remaining row \(a-2\), pair the \(B\)-columns in the same way.  This
tiles every remaining interior cell and leaves precisely the deleted corner.

\begin{figure}[htbp]
\centering
\begin{tikzpicture}[x=.48cm,y=.48cm,font=\scriptsize]
  \foreach \i in {0,...,6}{
    \foreach \j in {0,...,6}{
      \ifnum\i=0 \fill[orange!35] (\i,\j) rectangle ++(1,1);\fi
      \ifnum\i=6 \fill[orange!35] (\i,\j) rectangle ++(1,1);\fi
      \ifnum\j=0 \fill[orange!35] (\i,\j) rectangle ++(1,1);\fi
      \ifnum\j=6 \fill[orange!35] (\i,\j) rectangle ++(1,1);\fi
    }
  }
  \foreach \j in {1,...,5}{
    \foreach \i in {1,3}{
      \fill[blue!14] (\i,\j) rectangle ++(2,1);
      \draw[blue!70!black,very thick] (\i,\j) rectangle ++(2,1);
    }
  }
  \foreach \j in {1,3}{
    \fill[blue!14] (5,\j) rectangle ++(1,2);
    \draw[blue!70!black,very thick] (5,\j) rectangle ++(1,2);
  }
  \fill[white] (5,5) rectangle ++(1,1);
  \draw[red!75!black,very thick] (5,5) rectangle ++(1,1);
  \draw[red!75!black,thick] (5.18,5.18)--(5.82,5.82)
                              (5.18,5.82)--(5.82,5.18);
  \draw[step=1,black!65] (0,0) grid (7,7);
  \node[rotate=90] at (-.7,3.5) {\(A\)-turn coordinate};
  \node at (3.5,-.65) {\(B\)-turn coordinate};
  \node[align=left,anchor=west] at (7.45,5.8)
    {boundary singleton\\support};
  \draw[-{Latex[length=2mm]},orange!80!black] (7.35,5.55)--(6.55,5.55);
  \node[align=left,anchor=west] at (7.45,3.5)
    {interior domino\\support};
  \draw[-{Latex[length=2mm]},blue!70!black] (7.35,3.25)--(5.7,3.25);
  \node[align=left,anchor=west] at (7.45,1.2)
    {omitted majority-\\parity cell};
  \draw[-{Latex[length=2mm]},red!75!black] (7.35,1.35)--(5.65,5.35);
\end{tikzpicture}
\caption{The packing construction, shown for a \(7\times7\) turn grid.
Orange boundary cells and the four exceptional cycles have singleton support;
blue rectangles are pairwise disjoint two-turn supports.  The construction for
general odd \(a,b\) is the same row-first domino tiling with one majority-parity
interior cell omitted.}
\label{fig:packing-grid}
\end{figure}

\begin{theorem}[Symbolic packing bound]
\label{thm:packing-lift}
For odd \(a,b\ge3\),
\[
 \hsep(D(a,b))\ge \frac{(a+2)(b+2)-1}{2}.
\]
\end{theorem}

\begin{proof}
The four supports in \cref{eq:exception-requirements}, the boundary singleton
supports, and the domino supports in
\cref{eq:domino-support-a,eq:domino-support-b} are mutually disjoint: they
partition all Hamiltonian-cycle classes except \(R\) and the deleted interior
turn.  The all-rails cycle covers none of the chosen requirements.  Thus every
Hamiltonian cycle covers at most one chosen requirement.  Their number is
\[
 4+(2a+2b-4)+\frac{(a-2)(b-2)-1}{2}
 =\frac{(a+2)(b+2)-1}{2}.
\]
The packing half of Lemma~\ref{lem:certificate-principle} proves the claim.
\end{proof}

\subsection{A symbolic separating-family upper bound}

Define
\begin{align}
\label{eq:symbolic-family}
 \Ccal(a,b)={}&\{X_{0011},X_{1100},X_{0101},X_{1010}\}\notag\\
 &\cup\{T(i,j):i\in\{0,a-1\}\text{ or }j\in\{0,b-1\}\}\notag\\
 &\cup\{T(i,j):0<i<a-1,\ 0<j<b-1,\ i+j\text{ odd}\}.
\end{align}
There are \(2a+2b-4\) boundary turns.  The odd parity class of the odd by
odd interior grid has \(((a-2)(b-2)-1)/2\) cells, and hence
\begin{equation}
\label{eq:symbolic-family-size}
 |\Ccal(a,b)|=4+(2a+2b-4)+\frac{(a-2)(b-2)-1}{2}
 =\frac{(a+2)(b+2)-1}{2}.
\end{equation}

\begin{figure}[htbp]
\centering
\begin{tikzpicture}[x=.5cm,y=.5cm,font=\scriptsize]
  \foreach \i in {0,...,6}{
    \foreach \j in {0,...,6}{
      \ifnum\i=0 \fill[orange!35] (\i,\j) rectangle ++(1,1);\fi
      \ifnum\i=6 \fill[orange!35] (\i,\j) rectangle ++(1,1);\fi
      \ifnum\j=0 \fill[orange!35] (\i,\j) rectangle ++(1,1);\fi
      \ifnum\j=6 \fill[orange!35] (\i,\j) rectangle ++(1,1);\fi
      \ifnum\i>0\ifnum\i<6\ifnum\j>0\ifnum\j<6
        \pgfmathtruncatemacro{\parity}{mod(\i+\j,2)}
        \ifnum\parity=1 \fill[violet!28] (\i,\j) rectangle ++(1,1);\fi
      \fi\fi\fi\fi
    }
  }
  \draw[step=1,black!65] (0,0) grid (7,7);
  \foreach \i in {0,...,6}{
    \foreach \j in {0,...,6}{
      \ifnum\i=0 \node at (\i+.5,\j+.5) {\(\bullet\)};\fi
      \ifnum\i=6 \node at (\i+.5,\j+.5) {\(\bullet\)};\fi
      \ifnum\j=0 \node at (\i+.5,\j+.5) {\(\bullet\)};\fi
      \ifnum\j=6 \node at (\i+.5,\j+.5) {\(\bullet\)};\fi
      \ifnum\i>0\ifnum\i<6\ifnum\j>0\ifnum\j<6
        \pgfmathtruncatemacro{\parity}{mod(\i+\j,2)}
        \ifnum\parity=1 \node at (\i+.5,\j+.5) {\(\bullet\)};\fi
      \fi\fi\fi\fi
    }
  }
  \node[rotate=90] at (-.7,3.5) {\(i\)};
  \node at (3.5,-.65) {\(j\)};
  \node[align=left,anchor=west] at (7.5,5.7)
    {all boundary\\turns};
  \draw[-{Latex[length=2mm]},orange!80!black] (7.4,5.5)--(6.55,5.5);
  \node[align=left,anchor=west] at (7.5,3.3)
    {interior turns\\with \(i+j\) odd};
  \draw[-{Latex[length=2mm]},violet!80!black] (7.4,3.1)--(5.6,3.5);
  \node[align=left,anchor=west] at (7.5,1.1)
    {plus four\\\(X\)-cycles};
\end{tikzpicture}
\caption{Turn cycles selected by \(\Ccal(a,b)\), again illustrated on a
\(7\times7\) grid.  Every boundary turn is selected, as is one checkerboard
class of interior turns; the four connector-exception cycles are added
separately.}
\label{fig:separating-grid}
\end{figure}

We prove that this family separates by giving both directed witnesses for
every pair of edge types.  Let \(P\) be its selected turn cycles.  Write
\(R_i=\{T(i,j)\in P\}\), \(K_j=\{T(i,j)\in P\}\), and
\[
 N_A(p)=\{p-1,p\}\cap\{0,\ldots,a-1\},
 \qquad
 N_B(q)=\{q-1,q\}\cap\{0,\ldots,b-1\}.
\]
On the turn coordinates of an edge signature, every connector has all of
\(P\); an \(A\)-rail in cell \(i\) has \(P\setminus R_i\); an \(A\)-rung
\(p\) has precisely the selected turns whose first coordinate is in
\(N_A(p)\).  The corresponding statements for ladder \(B\) use columns.

For the four exception coordinates, abbreviate
\(B_0=X_{0011}\), \(B_1=X_{1100}\),
\(A_0=X_{0101}\), and \(A_1=X_{1010}\).  An \(A\)-rail contains
\(B_0,B_1\) and exactly one of \(A_0,A_1\), determined by its side and cell
parity.  An internal \(A\)-rung contains \(A_0,A_1\); an endpoint rung also
contains one of \(B_0,B_1\).  Interchanging \(A\) and \(B\) gives the other
ladder.  In connector order \cref{eq:connectors}, the four exception parts of
the connector signatures are
\begin{equation}
\label{eq:connector-signatures}
 \{B_1,A_1\},\quad\{B_1,A_0\},\quad
 \{B_0,A_1\},\quad\{B_0,A_0\}.
\end{equation}

\begin{theorem}[Symbolic separating family]
\label{thm:separating-family}
For odd \(a,b\ge3\), the edge signatures induced by \(\Ccal(a,b)\) are
pairwise incomparable.  Consequently,
\[
 \hsep(D(a,b))\le \frac{(a+2)(b+2)-1}{2}.
\]
\end{theorem}

\begin{proof}
We exhaust the eight edge-type combinations.  In every case we identify a
cycle in each of \(\sigma(e)\setminus\sigma(f)\) and
\(\sigma(f)\setminus\sigma(e)\).

\begin{enumerate}[label=\arabic*.,leftmargin=*]
\item \emph{Two connectors.}  The four sets in
\cref{eq:connector-signatures} are distinct and have size two.  Each of two
such sets therefore has an exception coordinate absent from the other.

\item \emph{A connector and a rail.}  A selected boundary turn in the rail's
omitted row or column belongs to the connector but not the rail.  A rail has
three exception coordinates---both opposite-ladder exceptions and one
alternating same-ladder exception---whereas a connector has only two; the
unused opposite-ladder exception belongs to the rail alone.

\item \emph{A connector and a rung.}  Since \(a,b\ge3\), there is a boundary
row or column outside the rung's one- or two-cell neighbourhood.  A boundary
turn there belongs only to the connector.  The rung contains both
same-ladder exceptions, while the connector contains exactly one of them, so
the other exception gives the reverse witness.

\item \emph{Two rails.}  The two sides of one cell have the same turn
coordinates but opposite alternating exception coordinates.  Rails in
different cells of one ladder are distinguished in both directions by
selected boundary turns in their two omitted rows or columns.  For rails on
different ladders, choose a boundary point in the omitted column but outside
the omitted row, and then one in the omitted row but outside the omitted
column.  Both are selected and give the two directions.

\item \emph{Two rungs of one ladder.}  If their neighbourhoods are not
nested, a selected boundary turn in each set difference gives both witnesses.
The only nested pairs are
\(N_A(0)\subset N_A(1)\) and
\(N_A(a)\subset N_A(a-1)\), and their \(B\)-analogues.  There the endpoint
rung has an additional exception coordinate, while a turn in the internal
rung's extra neighbouring cell gives the reverse witness.

\item \emph{Rungs of different ladders.}  Their two-element exception cores
are \(\{A_0,A_1\}\) and \(\{B_0,B_1\}\).  An endpoint rung can add at most
one coordinate from the other core, so each signature retains an exception
absent from the other.

\item \emph{A rail and a rung of the same ladder.}  The rail contains both
opposite-ladder exceptions and only one same-ladder exception.  The rung
contains both same-ladder exceptions and at most one opposite-ladder
exception.  Each therefore has an exception absent from the other.

\item \emph{A rail of one ladder and a rung of the other.}  Consider an
\(A\)-rail in row \(i\) and a \(B\)-rung with neighbourhood \(N_B(q)\).
Choose a boundary row different from \(i\) and a boundary column outside
\(N_B(q)\).  The corresponding selected turn belongs to the rail but not the
rung.  Conversely choose \(j\in N_B(q)\) so that \(T(i,j)\) is selected.
If \(i\) is a boundary row, every \(j\) works.  If \(i\) is interior, an
endpoint neighbourhood supplies a boundary column, while an internal
two-column neighbourhood contains exactly one column of the required parity.
Then \(T(i,j)\) belongs to the rung but is omitted by the rail.  The other
orientation is symmetric.
\end{enumerate}

The cases include endpoint rungs and both rail sides.  A square insertion
introduces only new rails and rungs with the same indexed incidences, so it
introduces no additional edge type or boundary case.  Thus all distinct edge
signatures are incomparable.  Apply
Proposition~\ref{prop:signature-antichain} and
\cref{eq:symbolic-family-size}.
\end{proof}

\subsection{The exact formula and its lift}

\begin{theorem}[Exact double-ladder formula]
\label{thm:exact-double-ladder}
For all odd integers \(a,b\ge3\),
\begin{equation}
\label{eq:exact-double-ladder}
 \boxed{\displaystyle
 \hsep(D(a,b))=\frac{(a+2)(b+2)-1}{2}.}
\end{equation}
\end{theorem}

\begin{proof}
The lower bound is Theorem~\ref{thm:packing-lift}; the upper bound is
Theorem~\ref{thm:separating-family}.  Their cardinalities agree.
\end{proof}

Both certificates lift locally under square insertion.  For example, under
\(b\mapsto b+2\), keep column zero and shift every old positive column by two.
The two new columns contribute four boundary cycles and one selected parity
cycle in each of the \(a-2\) interior rows, for a net increase of \(a+2\).
For the packing, the same strip contributes four boundary singletons and
\(a-2\) dominoes.  The \(a\mapsto a+2\) rule is symmetric and adds \(b+2\).
At the level of graph edges, turns at the expanded cell and two alternating
exceptions require the uniform local replacement inside the new square
gadget; all other cycles lift by replacing subdivided rails with their
three-edge paths.

\begin{figure}[htbp]
\centering
\begin{tikzpicture}[x=.43cm,y=.43cm,font=\scriptsize]
  \begin{scope}
    \foreach \i in {0,...,6}{
      \foreach \j in {0,...,4}{
        \ifnum\i=0 \fill[orange!28] (\j,\i) rectangle ++(1,1);\fi
        \ifnum\i=6 \fill[orange!28] (\j,\i) rectangle ++(1,1);\fi
        \ifnum\j=0 \fill[orange!28] (\j,\i) rectangle ++(1,1);\fi
        \ifnum\j=4 \fill[orange!28] (\j,\i) rectangle ++(1,1);\fi
        \ifnum\i>0\ifnum\i<6\ifnum\j>0\ifnum\j<4
          \pgfmathtruncatemacro{\parity}{mod(\i+\j,2)}
          \ifnum\parity=1 \fill[violet!24] (\j,\i) rectangle ++(1,1);\fi
        \fi\fi\fi\fi
      }
    }
    \draw[step=1,black!60] (0,0) grid (5,7);
    \node at (2.5,-.75) {\(a\times b\)};
  \end{scope}
  \draw[-{Latex[length=3mm]},thick] (5.7,3.5)--(7.0,3.5)
    node[midway,above] {insert two cells};
  \begin{scope}[xshift=8cm]
    \foreach \i in {0,...,6}{
      \foreach \j in {0,...,6}{
        \ifnum\i=0 \fill[orange!28] (\j,\i) rectangle ++(1,1);\fi
        \ifnum\i=6 \fill[orange!28] (\j,\i) rectangle ++(1,1);\fi
        \ifnum\j=0 \fill[orange!28] (\j,\i) rectangle ++(1,1);\fi
        \ifnum\j=6 \fill[orange!28] (\j,\i) rectangle ++(1,1);\fi
        \ifnum\i>0\ifnum\i<6\ifnum\j>0\ifnum\j<6
          \pgfmathtruncatemacro{\parity}{mod(\i+\j,2)}
          \ifnum\parity=1 \fill[violet!24] (\j,\i) rectangle ++(1,1);\fi
        \fi\fi\fi\fi
        \ifnum\j>0\ifnum\j<3
          \ifnum\i=0 \fill[green!45!black] (\j,\i) rectangle ++(1,1);\fi
          \ifnum\i=6 \fill[green!45!black] (\j,\i) rectangle ++(1,1);\fi
          \ifnum\i>0\ifnum\i<6
            \pgfmathtruncatemacro{\newparity}{mod(\i+\j,2)}
            \ifnum\newparity=1 \fill[green!45!black] (\j,\i) rectangle ++(1,1);\fi
          \fi\fi
        \fi\fi
      }
    }
    \draw[step=1,black!60] (0,0) grid (7,7);
    \draw[green!45!black,very thick] (1,0) rectangle (3,7);
    \node at (3.5,-.75) {\(a\times(b+2)\)};
    \node[align=center,anchor=west] at (7.5,3.5)
      {new selected turns:\\[2pt]\(4+(a-2)=a+2\)};
    \draw[-{Latex[length=2mm]},green!45!black,thick]
      (7.35,3.2)--(2.6,3.2);
  \end{scope}
\end{tikzpicture}
\caption{The separating-family lift for \(b\mapsto b+2\), illustrated with
\(a=7\).  Green cells are precisely the new selected turns: four boundary
turns and one parity turn in each interior row.  The packing lift has the same
increment, with the interior additions represented by domino supports.}
\label{fig:family-lift}
\end{figure}

The theorem concerns \(D(a,b)\) itself.  It does not assert that these graphs
maximize \(\hsep\) among all Barnette graphs of their orders.
